\subsection{مبانی نظری و فرمول‌بندی ریاضی}
فرض کنید $\mathbf{X} \in \mathbb{R}^{B \times N \times D}$ ماتریس حالت‌های پنهان ورودی باشد. ابتدا پروجکشن‌های پرس‌وجو، کلید و مقدار با استفاده از نگاشت‌های خطی مجهز به لورا تشکیل می‌شوند:
\begin{align}
\mathbf{Q} &= \operatorname{LoRA-Linear}_Q(\mathbf{X}) \in \mathbb{R}^{B \times H \times N \times d}, \\
\mathbf{K} &= \operatorname{LoRA-Linear}_K(\mathbf{X}) \in \mathbb{R}^{B \times H_{kv} \times N \times d}, \\
\mathbf{V} &= \operatorname{LoRA-Linear}_V(\mathbf{X}) \in \mathbb{R}^{B \times H_{kv} \times N \times d}
\end{align}

\paragraph{۱. انقباض خطی بعد دنباله (Sequence Contraction):}
ماژول مجهز به دو ماتریس پروجکشن دنباله یادگیرنده $\mathbf{P}_K, \mathbf{P}_V \in \mathbb{R}^{r \times N_{\text{max}}}$ است که بعد طول دنباله کلید و مقدار را از $N$ به $r$ منقبض می‌نمایند. برای طول دنباله جاری $N$، ماتریس برش‌خورده $\mathbf{P} \in \mathbb{R}^{r \times N}$ بر روی کلید و مقدار ضرب می‌شود:
\begin{align}
\mathbf{K}_{\text{proj}} &= \mathbf{P}_K \mathbf{K}_{\text{masked}} \in \mathbb{R}^{B \times H \times r \times d}, \\
\mathbf{V}_{\text{proj}} &= \mathbf{P}_V \mathbf{V}_{\text{masked}} \in \mathbb{R}^{B \times H \times r \times d}
\end{align}
پیش از اعمال این ضرب، به منظور ممانعت از نشت نویز توکن‌های پدینگ به داخل فضای فشرده زیرفضا، ماسک باینری معتبرسازی توکن‌ها اعمال می‌گردد:
\begin{equation}
\mathbf{K}_{\text{masked}} = \mathbf{K} \odot \mathbf{M}_{\text{valid}}, \quad \mathbf{V}_{\text{masked}} = \mathbf{V} \odot \mathbf{M}_{\text{valid}}
\end{equation}

\paragraph{۲. محاسبه ماتریس توجه نامتقارن و خروجی:}
ماتریس همبستگی توجه میان پرس‌وجو با ابعاد کامل $N$ و کلید فشرده با طول $r$ با پیچیدگی $\mathcal{O}(N \cdot r \cdot d)$ تشکیل می‌گردد:
\begin{equation}
\mathbf{S} = \frac{\mathbf{Q} \mathbf{K}_{\text{proj}}^\top}{\sqrt{d}} \in \mathbb{R}^{B \times H \times N \times r}
\end{equation}
\begin{equation}
\mathbf{A} = \operatorname{Softmax}(\mathbf{S}), \quad \mathbf{O}_{\text{heads}} = \mathbf{A} \mathbf{V}_{\text{proj}} \in \mathbb{R}^{B \times H \times N \times d}
\end{equation}
در نهایت، سرها ادغام شده و از لایه پروجکشن خروجی عبور داده می‌شوند:
\begin{equation}
\mathbf{Y} = \operatorname{LoRA-Linear}_O(\mathbf{O}_{\text{heads}}) \in \mathbb{R}^{B \times N \times D}
\end{equation}

که در آن $\sigma_i(\mathbf{W})$ بیانگر $i$-امین مقدار منفرد مرتب‌شده ماتریس به صورت نزولی است. نسبت حفظ انرژی طیفی $\eta_r = \frac{\sum_{i=1}^r \sigma_i^2}{\sum_{i=1}^{d} \sigma_i^2}$ تضمین می‌کند که با انتخاب آستانه $\tau \ge 0.90$، بیش از نود درصد از پراکندگی واریانس داده‌ها حفظ شده و افت سیگنال در انتشار گرادیان رو به عقب به حداقل ممکن کاهش می‌یابد.

\paragraph{تحلیل مجانبی و پیچیدگی اشتراک حافظه سخت‌افزاری}
از منظر پیچیدگی محاسباتی و انتقال داده‌ها در سلسومراتب حافظه (حافظه سراسری \lr{HBM} به حافظه محلی \lr{SRAM})، عملیات ضرب ماتریسی در این مکانیزم نیازمند تخصیص بافر موقت با ابعاد فاکتورگیری‌شده است. برای پردازش دنباله‌ای به طول $N$ با بعد مخفی $D$، بار محاسباتی هسته از فرم درجه دو به رابطه زیر تقلیل می‌یابد:
\begin{equation}
\label{eq:complexity_tlinformer_attention}
\mathcal{C}_{\text{tlinformer\_attention}} = \mathcal{O}\left( B \cdot N \cdot D \cdot r + B \cdot H \cdot N \cdot d_k \right) \ll \mathcal{O}\left( B \cdot H \cdot N^2 \cdot d_k \right)
\end{equation}
این بهبود مرتبه پیچیدگی، با کاهش چشمگیر حجم تبادل تانسورها در گذرگاه‌های ارتباطی کارت گرافیک، مانع از اشباع پهنای باند حافظه (پدیده \lr{Memory-Bound}) در فرآیند استنتاج طولانی‌مدت می‌گردد.
